%@ pages 4-16
%@ status translated
% Introduction (original section 1).
% Translated from the French of arXiv:math/9807049v3.
% The v2/v3 remarks are the authors' own and are kept in place.

\clearpage
\section{Introduction}
\label{intropage}

\subsection*{Descent: from modules to complexes}

To give an idea of the content of the present work, one may approach it, as a
first approximation, as a contribution to the theory of complexes of sheaves
of modules. To put it in perspective, we begin by summarising the theory of
sheaves of locally free modules (of fixed rank $r$):

There exists a stack $BG$, equipped with a universal module $U$ and a
universal bundle of which $U$ is the module of sections. This stack
``classifies'' locally free modules of rank $r$ on schemes, in the sense that
the stack of modules $Fib_{X,r}$ of such modules on a scheme $X$ is identified
with the stack of morphisms from $X$ to $BG$. The stack $BG$ is algebraic, and
if $X$ is projective then so is $Fib_{X,r}$. The stack $BG$ is an open subset
of (at least) two more general stacks $Coh$ and $Qcoh$, which classify
coherent and quasi-coherent sheaves respectively (we do not enter here into
the details, in particular those concerning the topology chosen). The fact
that $Qcoh$ is a stack means, more concretely, that descent data for a
quasi-coherent module (or for a morphism between two quasi-coherent modules)
are effective. One may observe that $Coh$ is not algebraic, which is a little
annoying.

We propose to generalise this theory to the case of complexes. The crucial
point is that we want to speak of gluing complexes together by means not of
isomorphisms but of \emph{quasi-isomorphisms} (compatible in a suitable
sense), so that this work belongs to a large extent to homotopy theory.

The stacks just mentioned are stacks of categories, and when one sets about
extending the preceding to the case of complexes, one must sooner or later
organise these complexes into categories equipped with suitable structures.

The solution that first comes to mind is to consider \emph{derived
categories}. One can indeed form, for example, the prestack $D_{Qcoh}$ which
to an affine scheme $Spec A$ associates the derived category $D_{Qcoh}(A)$ of
that of quasi-coherent $A$-modules. Asking whether this prestack is a stack
is the sophisticated way of asking whether one can glue complexes together by
means of ``compatible'' quasi-isomorphisms (i.e.\ satisfying a cocycle
condition that is otherwise rather technical). This question was recognised
very early on as misguided, the objects of these derived categories being
``essentially non-glueable in nature''\footnote{For the reader who is not
convinced that this route through derived categories is a dead end, let us
point out that the (full) subprestack $Parf^{\{0,1\}}$ of $D_{Qcoh}$ of
perfect complexes with cohomological support in $[0,1]$ is indeed a stack, but
that it is not locally algebraic (in the sense of Artin): in fact, as will be
seen later, the ``right'' object is a locally algebraic $2$-stack (in the
sense of \cite{geometricN}).} (\cite{SGA6} Exposé 0, p.~11).

\subsection*{Higher homotopies}

To formulate the problems of gluing complexes defined up to
quasi-isomorphism, one must consider higher homotopies, which complicates the
picture considerably. In ordinary topology, the prototype of such a problem
consists for example in gluing over $U_1 \cup U_2 \cup U_3 \cup U_4$ data of
the following kind: for $1 \leq i\leq 4$, $C_i$ is a complex on $U_i$; for
$1 \leq i\leq j \leq 4$, $f_{ij}$ is a homotopy equivalence between the
restrictions $C_{ij}$ and $C_{ji}$ of $C_i$ and $C_j$ to
$U_{ij}:=U_i \cap U_j$; for $1 \leq i\leq j \leq k\leq 4$, $h_{ijk}$ is a
homotopy equivalence between the restrictions of $f_{jk}\circ f_{ij}$ and
$f_{ik}$ to $U_{ijk}:=U_i \cap U_j \cap U_k$; and the $h_{ijk}$ must
moreover satisfy a compatibility condition on $U_1 \cap U_2 \cap U_3 \cap
U_4$, a condition whose very formulation is not immediate. One easily
conceives how, for more complex coverings (in the étale topology for
example), the combinatorics of this kind of descent data can become
intractable. The definition we give of descent data goes beyond (or avoids)
combinatorial considerations.

In the preceding situation the $f_{ij}$ are arrows between complexes, the
$h_{ijk}$ must be regarded as $2$-arrows, and the data on $U_1 \cap U_2 \cap
U_3 \cap U_4$ will be $3$-arrows. Concretely, $h_{ijk}$ for example is a
morphism of graded objects of degree $-1$ between $C_i$ and $C_k$ on
$U_{ijk}$, establishing a homotopy between $f_{jk}\circ f_{ij}$ and $f_{ik}$.
Likewise, the crucial information in a $3$-arrow is a morphism of graded
objects of degree $-2$ establishing a homotopy between morphisms of degree
$-1$, and so on. One therefore genuinely needs a notion of
\emph{$\infty$-category} as advocated by Grothendieck \cite{Grothendieck},
with $n$-arrows for every $n$ and the appropriate composition laws.

\subsection*{Simplicial or Segal categories}

In fact, a rudimentary form of $\infty$-category suited to our complexes has
been known for thirty years: it is the notion of a simplicial category.
Introduced in homotopy theory by Kan and Quillen (see \cite{Quillen}), this
notion was taken up by Dwyer and Kan \cite{DwyerKan1} \cite{DwyerKan2}
\cite{DwyerKan3}: in particular, to every category $M$ equipped with a
subcategory $W$, these authors associate a ``localised'' simplicial category
$L(M,W)$. The latter captures well the homotopy information concerning the
pair $(M,W)$, in the sense that if $M$ is a simplicial closed model category
(cmc) in Quillen's sense, and $W$ its subcategory of equivalences, then
$L(M,W)$ is equivalent to the simplicial category of cofibrant and fibrant
objects of $M$. This shows for example that the simplicial structure on a cmc
is unique up to equivalence, and indeed depends (again up to equivalence)
only on the subcategory of equivalences.

More particularly, in the situation of complexes, if $Ch$ denotes the
category of complexes and $qis$ the subcategory of weak equivalences, then
the simplicial category $L(Ch, qis)$ adequately represents the homotopy
theory of complexes.\footnote{This means that if $A^{\cdot}$ and $B^{\cdot}$
are complexes, then the simplicial set of morphisms in $L$ between
$A^{\cdot}$ and $B^{\cdot}$ is equivalent to the set obtained by applying the
Dold--Puppe construction to the truncated complex
$\tau ^{\leq 0}{\rm Hom}^{\cdot}(A^{\cdot}, B^{\cdot})$. One may note that a
simplicial category $L$ gives rise (following Grothendieck's idea
\cite{Grothendieck}) to an $\infty$-category $\Pi _{\infty}\circ L$ by
applying the ``Poincaré $\infty$-groupoid'' construction $\Pi _{\infty}$ to
each of the simplicial sets $Hom _L(x,y)$. In this $\infty$-category for
complexes, the $\infty$-category of morphisms between two complexes is
equivalent to $\gamma {\rm Hom}^{\cdot}(A^{\cdot}, B^{\cdot})$, where
$\gamma$ is the well-known construction which to a complex of abelian groups
$F^{\cdot}$ associates the strict $\infty$-category $\gamma (F^{\cdot})$
whose objects are the $x\in F^0$ with $d(x)=0$, whose morphisms between
$x,y$ are the $f\in F^{-1}$ with $d(f)=y-x$, whose $2$-morphisms between $f$
and $g$ are the $u\in F^{-2}$ with $d(u)=g-f$, and so on.}

For technical reasons, we systematically use the slightly more general
notion of a Segal category: a Segal category is a ``wide'' simplicial
category, in the sense that composition is not required to be strictly
associative. Every simplicial category is a Segal category and, conversely,
one shows that every Segal category is equivalent to a simplicial category
(\S7).

\subsection*{The problems posed}

Thus complexes of sheaves of modules organise themselves (for example over
the étale site) into presheaves of Segal categories. Among other things, one
can define ${\cal D}^+$ by ${\cal D}^+(Spec A)= L(Ch^+(A), qis)$, where
$Ch^+(A)$ is the category of complexes of $A$-modules bounded on the left and
$qis$ its subcategory of quasi-isomorphisms. One can now express as follows
the problems to which the present work is devoted:
\begin{itemize}
\item
identify the relevant descent problems concerning such presheaves;
\item
formulate these problems in the framework of a theory of (generalised)
stacks (in this introduction we shall say $\infty$-stacks) suited to Segal
categories; in this theory it will of course be necessary that the stacks are
the prestacks in which descent data glue together;
\item
identify a class of such presheaves, containing our presheaves of complexes,
in which these descent problems admit a solution.
\end{itemize}

Among the descent problems we want to treat there must figure the
``concrete'' gluing problems concerning complexes, of which we now give
examples. Let us point out however that we do not address these problems
directly in the present work---our aim being first of all to set up the
framework in which such questions should naturally be considered---and we
therefore formulate them only by way of motivation.

\begin{itemize}
\item
Let $k$ be a field, $X$ a regular $k$-scheme, and $\Ff$ a coherent sheaf on
$X$. It is known that $\Ff$ admits locally finite projective resolutions.
Can one glue these resolutions together into a global object?
\item
In \cite{SGA6} 0.4.4, for a morphism $f:X \rightarrow Y$ of schemes, one
constructs locally in $X$ a cotangent complex well defined up to
quasi-isomorphism. Can one glue these local constructions together?
\footnote{In loc.\ cit.\ it is considered that ``it is not possible to
proceed by simple gluing'', and the authors get around it differently.}

\item
Still in \cite{SGA6} (e.g.\ 0.4.2), one poses the problem (subsequently
solved by Deligne and Illusie, see \cite{Illusie} I.4.2) of defining the
exterior powers of a perfect complex, which we reformulate in our own way:
knowing that one can define, on bounded complexes of free modules of finite
type, a symmetric power functor $\Lambda^i$ which transforms
quasi-isomorphisms into quasi-isomorphisms, and knowing that perfect
complexes are those which are locally quasi-isomorphic to bounded complexes
of free sheaves of finite type, how is one to define the symmetric power of a
perfect complex?

\item
In the work of O'Brian, Toledo and Tong \cite{ObrianToledoTong1},
\cite{ObrianToledoTong2}, \cite{Toledo1}, \cite{Toledo2}, \cite{Toledo3}
devoted to another question arising from SGA 6, that of \emph{Riemann--Roch
formulas}, one finds \v{C}ech computations which are certainly an example of
a descent situation for complexes. A better general framework for these
computations could contribute to our understanding of Riemann--Roch formulas.

\item
There are certainly other potential applications than those concerning
complexes; see for example the work of Hinich \cite{Hinich1}, which treats a
descent problem for $1$-groupoids in relation to a corresponding descent
problem for presheaves of differential graded Lie algebras.
\end{itemize}

\subsection*{The \texorpdfstring{$\infty$}{infinity}-stacks}

Before presenting our view of descent data, we must explain our theory of
$\infty$-stacks. It has been understood above that we need---at a
minimum---a theory of stacks of simplicial (or Segal) categories, just as one
has a theory of stacks of sets (sheaves) and a theory of stacks of categories
(``classical'' stacks). And in a theory of stacks of simplicial categories,
the stack of morphisms between two objects (i.e.\ sections over an open
subset) will have to be a stack of simplicial sets. We would therefore like a
theory of $\Gamma$-stacks, where $\Gamma$ could indifferently be the category
of sets, or that of categories, or that of simplicial categories, or indeed
that of simplicial sets. A $\Gamma$-prestack on a site $\Xx$ is then a
presheaf on $\Xx$ with values in $\Gamma$ (we take strict presheaves for
simplicity, and because we hope that in the end it comes to the same thing).
And the (naïve) idea for expressing that a $\Gamma$-prestack $F$ is a
$\Gamma$-stack would consist in requiring that for every $X \in \Xx$ and
every sieve $\Bb$ covering $X$, $F(X)$ be the limit (in $\Gamma$) of the
restriction of $F$ to the category $\Bb$. If $\Gamma$ is the category of
sets one recovers of course the notion of sheaf, but if $\Gamma$ is the
category of categories, this no longer works. To rectify the situation and
fall back in this case on the usual stacks, it suffices to equip $\Gamma$
with its $2$-category structure and to require that $F(X)$ be the $2$-limit
(in $\Gamma$) of the restriction of $F$ to $\Bb$. To generalise this
situation, we shall therefore consider that $\Gamma$ is an $\infty$-category
\footnote{We do not take the trouble here to give a definition of
$\infty$-category: we simply regard Tamsamani's $n$-categories
\cite{Tamsamani}, as well as their variants introduced below, the Segal
$n$-categories, as $\infty$-categories.} in which one knows how to specify
homotopy limits (it is not necessary that they exist). One can then,
following a suggestion formulated by the second author in \cite{limits} 6.4,
say that a $\Gamma$-stack is a $\Gamma$-prestack $F$ such that for every
$X \in \Xx$ and every sieve $\Bb$ covering $X$, $F(X)$ is the limit (in
$\Gamma$) of the restriction of $F$ to the category $\Bb$.

\subsection*{The \texorpdfstring{$0$}{0}-stacks of Segal}

In the case where $\Gamma$ is the $\infty$-category (technically: the Segal
$1$-category) of simplicial sets, we confront successfully\footnote{Somewhat
tempered in the present version v3\ldots} this point of view with the
existing theory of simplicial presheaves, developed by K.~Brown
\cite{KBrown}, Illusie \cite{Illusie}, Joyal in a letter to Grothendieck
\cite{Joyal}, Jardine \cite{Jardine}, and Thomason \cite{Thomason}.\footnote{We
even have the $\infty$-category (a simplicial category here) of weak functors
from $\Xx ^o$ to $\Gamma$ thanks to the work of Cordier and Porter
\cite{CordierPorter}.} In this case one has at one's disposal the notion of
homotopy limit introduced by Bousfield and Kan \cite{BousfieldKan}, and our
$\Gamma$-stacks, which we call $0$-stacks of Segal, constitute a class of
simplicial presheaves already considered by Jardine (the presheaves flasque
with respect to every object of the site \cite{Jardine}) and, in the very
close framework of presheaves of spectra, by Thomason \cite{Thomason}. To
every simplicial presheaf (or Segal $0$-prestack) one knows how to associate
naturally a Segal $0$-stack having the same homotopy sheaves (and not
presheaves).

\subsection*{The \texorpdfstring{$n$}{n}-stacks of Segal}

It is therefore natural to extend this approach to the case of simplicial or
Segal categories. For this, we must equip the class $\Gamma$ of Segal
categories with an $\infty$-category structure (technically, it is a Segal
$2$-category) in which one can define the notion of (homotopy) limit. To
allow some room, we define the Segal $n$-categories outright (\S2): their
construction follows closely that of Tamsamani's $n$-categories
\cite{Tamsamani}. The Segal $0$-categories are the simplicial sets, and the
Segal $(n+1)$-categories are simplicial objects in the category of Segal
$n$-categories satisfying the (Segal) conditions which express that
composition, failing to be well defined and associative, is defined and
associative up to homotopy. The advantage of Segal $n$-categories over
simplicial sets lies in the fact that one can find in them non-invertible
$i$-morphisms, at least for $i \leq n$. Thanks to the notion of
\emph{internal} cmc (\S11), one shows that, up to set-theoretic
considerations, the Segal $n$-categories organise themselves into a Segal
$(n+1)$-category $nSeCAT$ (\S11), in which one knows how to define homotopy
limits (\S14). One can therefore introduce the ($\Gamma$-)stacks of Segal
$n$-categories, which we call Segal $n$-stacks. Technically, one begins by
giving a Jardine-style definition of Segal $n$-stacks (\S9), and it is only
in \S14 that one recovers the beautiful definition mentioned above. The
problem of descent for complexes then takes the following form: is the Segal
$1$-prestack of complexes a Segal $1$-stack?

\subsection*{Generalised descent data}

This formulation is not totally satisfactory, insofar as one does not see
clearly enough how it covers the concrete examples mentioned above.
Fortunately, we know how to introduce a notion of descent data (and also, of
course, of effective descent data) with values in a Segal $n$-prestack, which
plays the role expected of it. In fact, we explain here what is called
\emph{generalised descent data} (gdd), and we begin with the more familiar
case of simplicial presheaves: a gdd with values in the simplicial presheaf
$A$ on the site $\Xx$ is a morphism of simplicial presheaves
$\delta: D \rightarrow A$ where $D$ is \emph{contractible}, in the sense that
the Segal $0$-stack associated to $D$ is equivalent to the terminal object
$*_{\Xx}$. (Such a $D$ is essentially the same thing as a
\emph{hypercovering} in Verdier's sense \cite{SGA4b}.) Such data $\delta$ is
said to be \emph{effective} if it is equivalent to data
$\delta': D \rightarrow A$ which factors through $*_{\Xx}$. These definitions
extend without change to the case where $A$ is a Segal $n$-prestack, since
simplicial presheaves can also be regarded as Segal $n$-prestacks (\S2).

It is worth observing that the question of the effectivity of gdd in a Segal
$n$-prestack $A$ concerns only the \emph{interior} $A^{int}$ of $A$: the
interior of a Segal $n$-category $S$ is the simplicial set obtained from $S$
roughly by retaining only the morphisms (of all orders) which are invertible
up to homotopy, and this construction obviously extends to Segal
$n$-prestacks.

The preceding definition accounts well for the ``concrete'' examples
mentioned above. Thus, for the first example, one can take the (contractible)
simplicial presheaf $D$ such that $D(U)$ is the simplicial set\footnote{Here
as elsewhere in this work, we do not pay too much attention to questions
related to the difference between sets and classes.} of finite projective
resolutions of the restriction of $\Ff$ to $U$: the $1$-simplices are the
quasi-isomorphisms between resolutions (compatible with the augmentation)\ldots.

\subsection*{Characterisation of the \texorpdfstring{$n$}{n}-stacks of Segal by descent}

To explain in what way our notions of stacks are in turn suited to these
notions of descent data, we must introduce some notation. If $A$ is a Segal
$n$-category, the ($1$-)morphisms of $A$ organise themselves into a new Segal
$(n-1)$-category which we denote $Fl(A)$. If $x$ and $y$ are two objects of
$A$, we denote by $Fl(A)(x,y)$ the sub-Segal-$(n-1)$-category of $Fl(A)$ of
morphisms with source $x$ and target $y$. More generally, if $x$ and $y$ are
two objects of $Fl^{i-1}(A)$, we denote by $Fl^{i}(A)(x,y)$ the
sub-Segal-$(n-i)$-category of $Fl^i(A)$ of morphisms with source $x$ and
target $y$.\footnote{We adopt the convention that, for a given $i$, we
replace our Segal $n$-categories by equivalent Segal $n'$-categories, with
$n'>i$.} These notations extend to the case where $A$ is a Segal $n$-prestack
on a site $\Xx$: $Fl^i(A)$ is then a new Segal $(n-i)$-prestack on $\Xx$,
and if $x$ and $y$ are sections of $Fl^{i-1}(A)$ over an object $U$ of
$\Xx$, then $Fl^{i}(A)(x,y)$ is a Segal $(n-i)$-prestack on the site
$\Xx/U$. With these notations, one proves the following result:
\begin{theorem}
Let $A$ be a Segal $n$-prestack on the site $\Xx$. Suppose that the values of
$A$ are Segal $n$-categories (this condition is vacuous for $n=0$). Then the
following conditions are equivalent:
\newline
$(i)$ $A$ is a Segal $n$-stack;
\newline
$(ii)$
(a)\,\, for every $X\in \Xx$ and every $x,y\in A_0(X)$ the Segal
$n$-prestack $Fl(A)(x,y)$ on $\Xx /X$ is a Segal $n$-stack; and
\newline
(b) \,\, the descent data for $A$ are effective.
\newline
$(iii)$ (valid only for non-Segal $n$-stacks)\newline
(a')\,\, for every $i \ge 1$, for every $X\in \Xx$ and for all
$x,y\in Fl^{i-1}(A)(X)$ the descent data with values in the
$(n-i)$-prestack $Fl^{i}(A)(x,y)$ are effective; and
\newline
(b) \,\, the descent data with values in $A$ are effective.
\end{theorem}

This theorem formalises the idea (which emerges among other places from Giraud
\cite{GiraudThese} and Laumon--Moret-Bailly \cite{LaumonMB}) that a stack is a
prestack in which descent data are effective, and since the effectivity of
descent data concerns only the interior, it characterises our Segal
$n$-stacks in terms proper to the theory of simplicial presheaves.

\subsection*{Descent data and the coarse topology}

However, the very definition of (generalised) descent data formulated above
involved the notion of Segal $0$-stack, and it is therefore time to give our
definition of (true) descent data with values in a simplicial presheaf, since
it remains true that this case suffices. Let $\Xx$ be a site equipped with a
covering sieve $\Bb$ and with a simplicial presheaf $A$. Descent data on
$\Bb$ with values in $A$ is a diagram $\delta$ of simplicial presheaves:
$$
*_{\Bb} \leftarrow D \rightarrow A
$$
where the left-hand arrow is what we call an \emph{object-by-object
equivalence}, that is to say it induces a weak equivalence of simplicial sets
on each object $U \in \Xx$ (here, of course, only the objects of $\Bb$ are
really concerned). This implies that $D$ is contractible, and one is indeed
in the presence of descent data in the generalised sense given above. One may
observe that giving the diagram $\delta$ is equivalent to giving descent data
(generalised or not, it is the same) with values in the restriction of $A$ to
the category $\Bb$ equipped with the coarse topology. This is a thoroughly
typical example of the way the coarse topology intervenes systematically
throughout our work. Another typical example is the way we describe in two
stages the model category structure for our prestacks, by Bousfield
localisation from the structure corresponding to the coarse topology (\S6).

Descent data is said to be effective if it is equivalent to descent data
which factors through $*_{\Xx}$. And when in the theorem above we say that
the descent data with values in $A$ are effective, we mean that for every
object $X$ of $\Xx$ and every sieve $\Bb$ covering $\Xx/X$, every descent
datum on $\Bb$ with values in $A|_{\Xx/X}$ is effective.\ednote{The original
prints $A*{|\Xx/X}$; we read this as the restriction of $A$ to $\Xx/X$.}

\subsection*{The two stages of descent and Grothendieck's constructions}

The problem of the effectivity of descent data in a Segal $n$-prestack $A$ on
a site $\Xx$ (admitting a terminal object) decomposes naturally into two
stages. Indeed, descent data $\delta$ as above fits into a diagram:
$$
*_{\Xx} \leftarrow *_{\Bb} \leftarrow D \rightarrow A,
$$
Since the point is to factor, up to equivalence, the given morphism
$D \rightarrow A$ through $*_{\Xx}$, one may first look for such a
factorisation through $*_{\Bb}$. If therefore one calls strict the descent
data with $D=*_{\Bb}$, one can formulate on the one hand the problem of the
strictification of descent data, which consists in showing that every descent
datum is equivalent to a strict descent datum, and on the other hand that of
the effectivity of strict descent data.

Technically, things present themselves a little differently. To show that a
Segal prestack $A$ on $\Xx$ is a Segal stack, one uses the definition by
limits. One must therefore compare $A(X)$ with the homotopy limit mentioned
above, say $\Gamma(\Bb, A)$. To identify this limit, one follows the path
traced in \cite{SGA1}: one constructs a Segal category $\int _{\Bb} A$ over
$\Bb^o$ in the hope of identifying $\Gamma(\Bb, A)$ with a category of
sections of $\int _{\Bb} A \rightarrow \Bb^o$. In fact, one finds (\S16) that
$\Gamma(\Bb, A)$ is equivalent to the Segal category of ``balanced'' sections
(or eq-sections), not of $\int _{\Bb} A \rightarrow \Bb^o$ but of its fibrant
replacement $\int '_{\Bb} A \rightarrow \Bb^o$.\ednote{The French is
\emph{sections équilibrées}, abbreviated \emph{eq-sections}; this is the
authors' own term, not an established English one. The \emph{eq} abbreviates
\emph{équilibrée} (``balanced''), not ``equi-free'' or ``equivariant''. In
\S16 an eq-section is defined as a section $\sigma$ for which the maps
$\sigma^r(f)$ are equivalences. We keep the authors' abbreviation
\emph{eq-section} throughout and use ``balanced'' only here, at the first
occurrence.} This replacement refers to a
cmc structure for Segal categories (\S2). One therefore wants to show that a
composite morphism
$$
A(X) \rightarrow
Sect^{eq}(\int _{\Bb} A) \rightarrow
Sect^{eq}(\int' _{\Bb} A)
$$
is an equivalence. One may regard the right-hand term as a category of
descent data (on the right?) and the middle one as the subcategory of strict
descent data (we have not sought to establish a correspondence between these
new notions of descent data and the old ones), and the problem of effectivity
decomposes again into two stages.

These two stages will be treated each in turn in our \S18 and \S19. To remain
strictly in conformity with what happens in the paper, one should say that
this situation of ``strictification'' of descent data takes place not over
$\Bb$ but over the category $\Delta$ via a functor
$\rho (U):\Delta \rightarrow \Bb$, in the case where $\Bb$ is the sieve
generated by an object $U$ covering $X$. In a way $\Delta$ approximates
$\Bb$ via the functor $\rho (U)$, and for technical reasons it is more
convenient for us to work over $\Delta$.

It so happens that the problem of the effectivity of strict descent data has
already been solved for our complexes of modules on a ringed topos in
Exposé Vbis of SGA 4 II \cite{SGA4} by Saint-Donat (after ``succinct notes''
of Deligne). In \S19 we give a more general result for this ``strict'' stage.

\subsection*{Quillen presheaves}

We have thus set up this formalism of Segal stacks which constitutes a
natural framework for relative homotopy theory (over a site), in which we can
reformulate the problem of descent: to find conditions for a Segal prestack
to be a (Segal) stack. As Jardine writes in \cite{Jardine}, the organisational
structure for homotopy theory is that of a closed model category (cmc) in
Quillen's sense \cite{Quillen}. The theory of Segal stacks that we have set
up is the natural framework of relative homotopy theory (over a site), and
the corresponding organisational structure is quite naturally a structure of
presheaf of cmc's: we call Quillen presheaf any presheaf of cmc's in which
the restrictions are what are called (left) Quillen functors. If ${\bf M}$ is
such a Quillen presheaf on a site $\Xx$, one can associate to it the Segal
$1$-prestack $L({\bf M})$ on $\Xx$ obtained by taking for $L({\bf M})(X)$ the
simplicial category obtained by inverting (in the sense of Dwyer--Kan) the
equivalences in ${\bf M}(X)$. One therefore seeks sufficient conditions for
the Segal prestack thus associated to a Quillen presheaf to be a Segal stack.

\subsection*{A descent theorem}

We therefore propose a theorem (\ref {dansleschamps}) giving sufficient
conditions for the prestack $L({\bf M})$ to be a Segal stack. It is of course
starting from the case of complexes that we found such a general formulation.
The conditions we impose are of three kinds.

First, we restrict ourselves to sites where, roughly speaking, one can reduce
to considering only descent data defined on coverings with a single object
(and satisfying the appropriate compatibility conditions on the \v{C}ech
complex of the covering).

The main restriction that we impose on the Quillen presheaf is the existence,
in the values taken (which are cmc's), of arbitrary small limits and
colimits. Thus for example our theorem does not concern complexes of sheaves
of finite type. We also impose on it other more or less natural conditions,
namely: the fact that, in the values taken, all objects are cofibrant; the
existence of functorial factorisations; compatibility with fibre products;
the local character of equivalences; and compatibility with disjoint unions.

Finally we impose a very technical condition allowing the use of Quillen's
so-called ``small object'' argument for the construction of an auxiliary cmc.

The most delicate part of the proof consists in the strictification of
descent data (\S18).

\noindent
--------
\newline
In version 2 of the paper we had to add hypothesis (4), that the restriction
morphisms for $L(\mm )$ are exact on both sides (and not only on the side
corresponding to the fact that $\mm$ is left Quillen). Indeed, the first
version contained an error in one line of the proof, where we commuted a
restriction morphism with a limit; we add this commutation as a hypothesis
since it is not a consequence of the others. This condition seems to be
primordial; fortunately it is verified in the applications in \S20 and \S21.

\subsection*{Stacks of stacks}

It was in studying our descent problem for complexes that we realised there
was an analogous result for $n$-stacks themselves: the $n+1$-prestack
$n\underline{CHAMP}$ defined by
$$
n\underline{CHAMP}(X):= nCHAMP (\Xx /X)
$$
is an $n$-stack (see Theorems \ref{descente} and \ref{descenteNdS}). This is a
``gluing'' result for $n$-stacks. For example, given stacks $\Ff _U$ and
$\Ff _V$ on two open subsets $U,V$ of a topological space, with an
equivalence
$$
\Ff _U |_{U\cap V} \cong
\Ff _V |_{U\cap V}
$$
there then exists a stack $\Ff $ on $U\cup V$ with $\Ff |_U \cong \Ff _U$ and
$\Ff |_V \cong \Ff _V$. One can regard \ref{descente} and \ref{descenteNdS}
as generalisations to $n$-stacks of the classical result which says that the
$1$-prestack which to each open subset $U$ associates the category of sheaves
of sets (i.e.\ $0$-stacks) on $U$ is a $1$-stack. The case $n=1$ of this
generalisation is already in SGA 1 \cite{SGA1} \cite{GiraudThese}, where it is
shown that the $2$-prestack $1\underline{CHAMP}$ of $1$-stacks is a
$2$-stack.

This result for $n=2$ (which says that $2\underline{CHAMP}$ is a $3$-stack)
was used, without proof, by Breen \cite{Breen}. One may say that our result
generalises approximately half of what Breen does in \cite{Breen} (the other
half being the expression of descent data in terms of cocycles, something we
have not managed to do in a completely satisfactory way for arbitrary $n$).

In fact, at the beginning of this work, the only thing we could prove about
complexes was that the prestack of complexes of $\qq$-vector spaces, of fixed
length $n$, was an $n+1$-stack. To prove this, we had used the result saying
that the $n+1$-prestack of $n$-stacks is a stack: with this and an elementary
version of the notion of ``spectrum'' (consisting just in making a translation
so that everything happens between degrees $N$ and $N+n$ with $N\geq n+2$)
one could obtain (thanks to the fact that stable rational homotopy is
trivial) the fact that the complexes of $\qq$-vector spaces of length $n$
form an $n+1$-stack.

It is only with the contribution of state-of-the-art techniques in cmc theory
(\cite{DwyerKan1} \cite{DwyerKan2} \cite{DwyerKan3} \cite{DwyerKanDiags},
\cite{Hirschhorn},\cite{DHK}, \cite{JardineGoerssBook}, \cite{HoveyMonMod}
\cite{HoveyBook}) that we have been able to formulate the more general result
\ref{dansleschamps}, which allows us to treat the case of complexes of
$\Oo$-modules.

\subsection*{Algebraicity}

The more detailed study of stacks of complexes on sites of algebraic geometry
will be the subject of another work. We simply state in \S21 our main
algebraicity result.

\subsection*{History, related work and techniques used}

The notion of $1$-stack is today well established in algebraic geometry,
notably since the work of Artin \cite{ArtinInventiones}, of
Deligne--Mumford \cite{Deligne-Mumford}, of Laumon--Moret-Bailly
\cite{LaumonMB}, extended by numerous current works. A similar notion plays
an important role in differential geometry (e.g.\ in the geometry of
foliations), see for example \cite{Ehresmann} \cite{Morita} \cite{Pradines}
\cite{Haefliger} \cite{Haefliger2} \cite{Molino} \cite{Moerdijk}.
%
\footnote{One may cite Moerdijk (\cite{Moerdijk} Example 1.1 (f)):
``Effective etale groupoids are well-known to be basically equivalent to
pseudogroups\ldots'' followed by a reference to Molino \cite{Molino}. In view
of this remark, the work of Ehresmann \emph{et al} concerns essentially
étale groupoids, i.e.\ Deligne--Mumford $1$-stacks in the differential
setting. It is not clear to our eyes whether this connection was known to
Ehresmann at the time; in any case, this point of view is already present in
the papers of Pradines \cite{Pradines}; and according to Pradines (oral
communication to the second author), this point of view was presented by
Ehresmann in his talk at the international congress of 1956.}
%
Building on the notion of stack, Giraud developed non-abelian cohomology in
low degrees \cite{Giraud}.

A precursor work on descent for the $\infty$-stack of complexes is that on
``cohomological descent'' by B.~Saint-Donat and Deligne in SGA 4
\cite{SGA4b}.

In \cite{Grothendieck}, A.~Grothendieck proposes as a generalisation the
notion of \emph{$n$-stack} or even of \emph{$\infty$-stack}, without however
giving a precise definition. In fact, one may say that the theory of
$\infty$-stacks of groupoids is already available at the time of
\cite{Grothendieck} by way of the theory of simplicial presheaves, thanks to
the work of K.~Brown \cite{KBrown}, Illusie \cite{Illusie}, Joyal in a letter
to Grothendieck \cite{Joyal}, Jardine \cite{Jardine}, and Thomason
\cite{Thomason}. The way in question passes through the equivalence between
the theory of $\infty$-groupoids and that of topological spaces or simplicial
sets---an equivalence for which only the definition of $\infty$-groupoid was
missing!

More recently, the theory of $n$-categories for small values of $n$ has been
recognised as important in a wide range of situations linked for example to
new knot invariants, to quantum groups, to quantum field theory, and even to
computer science\ldots. This has led to a better understanding in particular
of the role played by the non-associativity of composition. The first paper
of Baez--Dolan \cite{BaezDolan} provides an interesting summary of these
developments.

Tamsamani's thesis \cite{Tamsamani} provides a definition of $n$-category, as
well as the equivalence between $n$-groupoids and $n$-truncated topological
spaces, for every $n$. This definition makes it possible to approach the
theory of $n$-stacks, not necessarily of groupoids, using the techniques due
to Quillen which underlie the work of Brown, Illusie, Joyal, Jardine and
Thomason. Let us also point out that other definitions of $n$-category have
been proposed, see Baez--Dolan (\cite{BaezDolanLetter} \cite{BaezDolanIII},
which was independent of \cite{Tamsamani} and simultaneous) or Batanin
(\cite{Batanin}, a little later). For the moment a comparison between these
different definitions is missing, and as explained above, we use a slight
variant of Tamsamani's definition.

Our basic technique is that of cmc's, and our work has been deeply influenced
by recent work in cmc theory, that of Hirschhorn \cite{Hirschhorn}, of
Dwyer--Hirschhorn--Kan \cite{DHK}, also that of Goerss--Jardine
\cite{JardineGoerssBook} and of Hovey \cite{HoveyMonMod} \cite{HoveyBook}, as
well as by the less recent work of Dwyer--Kan on localisation \cite{DwyerKan1}
\cite{DwyerKan2} \cite{DwyerKan3} \cite{DwyerKanDiags}.

\subsection*{Acknowledgements}

We thank C.~Walter, N.~Mestrano, Z.~Tamsamani, P.~Hirschhorn, B.~Toën,
R.~Brown, A.~Bruguières, T.~Goodwillie, L.~Katzarkov, M.~Kontsevich,
G.~Maltsiniotis, H.~Miller, S.~Mochizuki, T.~Pantev, J.~Pradines, C.~Rezk,
J.~Tapia, C.~Teleman, D.~Toledo for discussions, correspondence, courses,
theses and books which have influenced the present work. The second author
thanks the University of California at Irvine for its hospitality at the end
of the preparation of this paper.

We thank our families for their support during this work.
